\section*{2. The Integers}
\addcontentsline{toc}{section}{2. The Integers}

\begin{exercise}{2.1}
\begin{enumerate}[label=(\alph*)]
\item Find the greatest common divisor of \(34\) and \(21\), using Euclid's Algorithm. Then express their gcd as a linear combination of \(34\) and \(21\).
\item Now do the same for \(2424\) and \(772\).
\item Do the same for \(2007\) and \(203\).
\item Do the same for \(3604\) and \(4770\).
\end{enumerate}
\end{exercise}
\begin{proof}
    Tedious.
\end{proof}

\begin{exercise}{2.2}
\begin{enumerate}[label=(\alph*)]
\item Prove that \(\gcd(a,b)\) divides \(a-b\). This sometimes provides a short cut in finding \(\gcd\)s.
\item Use this to find \(\gcd(1962, 1965)\).
\item Now find \(\gcd(1961, 1965)\).
\end{enumerate}
\end{exercise}
\begin{proof}
\begin{enumerate}[label=(\alph*)]
\item By definition, \(\gcd(a,b)|a\) and \(\gcd(a,b)|b\). So \(a = \alpha\gcd(a,b)\) and \(b = \beta\gcd(a,b)\) for some \(\alpha,\beta\).
Hence, \(a-b = \alpha\gcd(a,b) - \beta\gcd(a,b) = \gcd(a,b)(\alpha - \beta)\). Therefore, \(\gcd(a,b) | a-b\).
\item \(1965 - 1962 = 3\). The only factors of \(3\) are \(1, 3\). \(3 | 1962\) and \(3 | 1965\), so \(\gcd(1962, 1965) = 3\).
\item \(1965 - 1961 = 4\). The factors of \(4\) are \(1, 2, 4\). Clearly \(4\) and \(2\) cannot be divisors of \(1961\) nor \(1965\). Hence, \(\gcd(1961, 1965) = 1\).
\end{enumerate}
\end{proof}

\begin{exercise}{2.3}
    Prove that the set of all linear combinations of \(a\) and \(b\) are precisely the multiples of \(\gcd(a,b)\).
\end{exercise}
\begin{proof}
    \[S := \{ax+by\ :\ \text{for arbitrary }x,y \in \integers\}.\]
    Let \(k \in S\) be an arbitrary element.
    We know \(\gcd(a,b) \mid a\) and \(\gcd(a,b) \mid b\), therefore there exists \(\alpha, \beta\) such that \(a = \alpha\gcd(a,b)\) and \(b = \beta\gcd(a,b)\).
    We also know that \(k = ax+by\) for some \(x,y \in \integers\). With substitution, we see that \[k = (\alpha\gcd(a,b))x + (\beta\gcd(a,b))y = \gcd(a,b)(\alpha x+\beta y).\]
    Now let \(k \in \integers\) be arbitrary, then \(k\gcd(a,b)\) is an arbitrary multiple of \(\gcd(a,b)\).
    From the GCD identity, \(\gcd(a,b) = ax+by\) for some \(x,y \in \integers\). Thus, \[k\gcd(a,b) = k(ax+by) = akx + bky \in S.\]
\end{proof}

\begin{exercise}{2.4}
    Two numbers are said to be \textbf{relatively prime} if their \(\gcd\) is \(1\). Prove that \(a\) and \(b\) are relatively prime
    if and only if every integer can be written as a linear combination of \(a\) and \(b\).
\end{exercise}
\begin{proof}
    Let \(a,b\) be arbitrary relatively prime integers. By definition, \(\gcd(a,b) = 1\). It follows
    from the GCD identity that there exists \(x,y \in \integers\) such that \(\gcd(a,b) = ax+by = 1\).
    Let \(k\) be an arbitrary integer, then \(k(ax+by) = a(kx)+b(ky) = k\) is a linear combination of \(a\) and \(b\) that equals \(k\).
    Now suppose that every integer can be written as a linear combination of \(a\) and \(b\). In particular, there exists \(x,y \in \integers\) such that
    \(ax+by = 1\). From Corollary \(2.5\), it follows immediately that \(ax+by = 1 = \gcd(a,b)\). Hence, \(a\) and \(b\) are relatively prime.
\end{proof}

\begin{exercise}{2.5}
    Prove Theorem 2.6. That is, use induction to prove that if the prime \(p\) divides \(a_1 a_2 \cdots a_n\), then \(p\) divides \(a_i\), for some \(i\).
\end{exercise}
\begin{proof}
    \footnotemark We induct on \(n\). If \(n = 1\), then the statement is trivial. And \(n = 2\) is true because \(p\) is prime.
    Now suppose that the statement is true for all \(k < n\) where \(n > 2\). In particular, the statement is true for \(k = n-1\).
    Suppose that \(p\) divides \(a_1 a_2 \cdots a_n\), from \(p\)'s prime properties, we know that \(p \mid a_n\) or \(p \mid a_1 a_2 \cdots a_{n-1}\).
    Suppose the former, then the statement holds. Suppose the latter, from our induction hypothesis, we know that \(p \mid a_i\) for some \(1 \leq i \leq n-1\).
    This closes the induction.
\end{proof}
\footnotetext{Starting from now, I'll refrain from using overly formal induction proofs that involve set definitions.}

\begin{exercise}{2.6}
    Suppose that \(a\) and \(b\) are positive integers. If \(a+b\) is prime, prove that \(\gcd(a,b) = 1\).
\end{exercise}
\begin{proof}
    Suppose that \(a + b\) is prime. Assume, for the sake of contradiction, that \(\gcd(a,b) > 1\).
    By definition, there exists \(\alpha, \beta \in \integers\) such that
    \(a = \alpha\gcd(a,b)\) and \(b = \beta\gcd(a,b)\). Because \(a > 0\), \(b > 0\) and \(\gcd(a,b) > 0\), we deduce
    that \(\alpha > 0\) and \(\beta > 0\). By substitution, \(a + b = \gcd(a,b)(\alpha + \beta)\).
    From Theorem \(2.7\), we know that \(a+b\) must be irreducible. But \(\gcd(a,b) > 1\) and \(\alpha + \beta > 1\).
    This contradiction shows that \(\gcd(a,b) = 1\), as desired.
\end{proof}